\documentclass{ximera}

\input{../preamble}


\title{Our Problems are Multiplying}
\author{Vic Ferdinand, Betsy McNeal, Jenny Sheldon}

\begin{document}
\begin{abstract}
    Let's try multiplying some integers!
\end{abstract}
\maketitle



Here are some multiplication problems we'll be working with throughout the activity.

\begin{enumerate}
    \item $5 \times 2$
    \item $15 \times (-3)$
    \item $(-4) \times 9$
    \item $(-6) \times (-12)$
\end{enumerate}


\begin{problem}
Use red and black chips to model each of the expressions above. How do you see each of the numbers? How do you see the operation in your picture?

\end{problem}



\begin{problem}
Use a number line to model each of the expressions above. How do you see each of the numbers? How do you see the operation in your picture?  
\end{problem}

\begin{problem}
Use checks and bills to model each of the expressions above. How do you see each of the numbers? How do you see the operation in your picture?  
\end{problem}




\newpage
\begin{instructorNotes}
This activity is the third of three that investigate operations with negative numbers.  The first is ``Where Have the Integers Been?'' and the second is ``Integers and Subtraction''.  There is an accompanying reading (\url{https://ximera.osu.edu/m4t/elementaryTeachersOne}) on integers that lays out some of the more complicated framework, and the instructor notes for ``Where Have the Integers Been?'' contain background information about this unit.

The main purpose of this final activity is to talk about issues with multiplication, as well as clear up as much confusion and as many questions as we can.  Our goal is to have students be familiar with number lines and checks and bills by the end of the activity.

This activity's questions are intentionally similar to those on the other handouts about integers, so that we can quickly remind our students of the previous activity, and then select the problems where they need the most practice.  We do not often discuss all of the problems on this page.

With multiplication, we take extreme care with the checks and bills method.  Again, we attend carefully to the question being asked, to ensure that it is answered by the appropriate expression.  For instance, in the story, ``I receive 3 checks, each for 5 dollars.  What is my balance now?''.  Here, ``what is my balance?'' is not a good question for this story problem, but we could adjust it by explicitly starting with a zero balance.  In the same vein, the question ``how much has my balance gone down?'' should have a positive answer, even though the question suggests it should be negative.  We find that students need a lot of practice evaluating stories for correctness.

Stories that correspond with multiplication problems with a negative number in the first factor can still feel contrived to students.  If this comes up, we remind them again that all of these contexts have issues, and that none is perfect!  To help with these issues, you may want to think about starting with a zero balance and then adding or taking (receiving or sending) from this zero with multiplication stories.  

As we discussed with subtraction, we also would like to use this activity to help students see the reasoning behind the usual rule ``a negative times a negative gives a positive''.  Both the number line context as well as the checks and bills context (sending bills that will be paid off with received checks) can be used to help make this idea clear.  You might also consider an arithmetic expression in addition or subtraction to correspond with the story (e.g., (-3) times (-5) is the same as 0 - (-5) - (-5) - (-5)).

{\bf Suggested Timing:} This activity takes us an entire class period.  We often do a 5-minute reminder of the previous work, then give students 5 minutes on number lines, and then 15 minutes to discuss the number line context.  We then give 5-10 minutes for checks and bills, and the remainder of the time to discussion.  We often swap the order of these contexts as well, or just talk about checks and bills if the students are doing well with number lines.  If there are enough lingering questions, we sometimes spend a small amount of the next class clearing those up.
\end{instructorNotes}





\end{document}
